\section{Estado qutrit, purificaci\'on termocampo doble y d\'eficit informacional}

Para un enlace con peso \(w>0\), sean
\begin{equation}
\label{rm:eq:modular-link-state}
Z(w)=1+\ee^{-w}+\ee^{-2w},
\qquad
\rho(w)=\frac1{Z(w)}\sum_{n=0}^2\ee^{-wn}|n\rangle\langle n|.
\end{equation}
Los pesos de las dos familias son
\begin{align}
w_{90}&=90\Dmod
=0.0100077345053979592447310951874\ldots,
\label{rm:eq:modular-w90}\\
w_{120}&=120\Dmod
=0.0133436460071972789929747935832\ldots,
\qquad \frac{w_{120}}{w_{90}}=\frac43.
\label{rm:eq:modular-w120}
\end{align}
El estado finito de borde es
\begin{equation}
\label{rm:eq:modular-product-state}
\rho_A=\rho(w_{90})^{\otimes18}
\otimes\rho(w_{120})^{\otimes18}.
\end{equation}
Su logaritmo reproduce \cref{rm:eq:modular-hamiltonian}, con
\(c=18\log Z(w_{90})+18\log Z(w_{120})\).

\begin{theorem}[Purificaci\'on termocampo doble]
\label{rm:thm:modular-tfd}
El vector
\begin{equation}
\label{rm:eq:modular-tfd-link}
|\operatorname{TFD}(w)\rangle
=\frac1{\sqrt{Z(w)}}\sum_{n=0}^2
\ee^{-wn/2}|n\rangle_A\otimes|n\rangle_{\bar A}
\end{equation}
purifica \(\rho(w)\). El producto de dieciocho copias para cada uno de los
pesos \(w_{90},w_{120}\) purifica \(\rho_A\); en el sistema duplicado,
\(S(\rho_A)\) es exactamente entrop\'ia de entrelazamiento.
\end{theorem}

\begin{proof}
Al trazar el segundo factor, la ortogonalidad de \(|n\rangle_{\bar A}\)
elimina los t\'erminos cruzados y deja
\[
\operatorname{Tr}_{\bar A}
|\operatorname{TFD}(w)\rangle\langle\operatorname{TFD}(w)|
=\rho(w).
\]
La traza parcial conmuta con productos tensoriales, lo que prueba la
afirmaci\'on para los treinta y seis enlaces.
\end{proof}

Para un enlace,
\begin{equation}
\label{rm:eq:modular-link-entropy}
S(w)=\log Z(w)
+w\frac{\ee^{-w}+2\ee^{-2w}}{Z(w)}.
\end{equation}
En consecuencia,
\begin{equation}
\label{rm:eq:modular-entropy-value}
S(\rho_A)=18S(w_{90})+18S(w_{120})
=39.54837320881807994957319730\ldots\ \text{nats}.
\end{equation}
El estado m\'aximamente mezclado de treinta y seis qutrits tiene entrop\'ia
\(36\log3\). Definimos el defecto orientado
\begin{equation}
\label{rm:eq:modular-oriented-information}
\boxed{
I_{\rm or}=36\log3-S(\rho_A)
=0.001669183233868940655631225913\ldots\ \text{nats}.}
\end{equation}
Su positividad es consecuencia de la maximalidad entr\'opica del estado
uniforme. La cercan\'ia de \(I_{\rm or}\) a \(\gammaH a_0\) no constituye
una identidad y no se usa para seleccionar ninguno de los dos valores.
